\documentclass[11pt]{article}
\usepackage[a4paper,margin=20mm]{geometry}
\usepackage{amsmath,amssymb,amsthm,microtype}
\usepackage[colorlinks=true,linkcolor=blue,urlcolor=blue]{hyperref}
\newtheorem{theorem}{Theorem}
\title{Exact circuits for MPUs with a fixed physical-entry alphabet}
\author{TNLean contributors}
\date{3 October 2026}
\begin{document}
\maketitle

An exact circuit bound polynomial jointly in chain length and
operator Schmidt rank is available when the entries of the full
physical matrix belong to a fixed finite alphabet. The alphabet
provides two gaps: a nonzero probability gap for recovering support,
and a separation gap for recovering exact coefficient labels.
Two coherently prepared isometries then have overlap equal to a
constant multiple of the target unitary. Exact uniform-amplitude
amplification removes that constant. This result allows coherent
mixing of a bounded number of basis outputs and extends the
permutation case. It is a mathematical proof and is not asserted
to be formalized in Lean.

\section{Exact joint polynomial synthesis for a fixed physical-entry alphabet}

Fix $d\ge2$ and a finite nonempty set
$\Xi\subset\mathbb C\setminus\{0\}$. Define
\[
 \eta=\min_{\alpha\in\Xi}|\alpha|>0,\qquad
 \delta=\min_{\substack{\alpha,\beta\in\Xi\cup\{0\}\\
                              \alpha\ne\beta}}|\alpha-\beta|>0,
 \qquad K=\lceil\eta^{-2}\rceil.
\]
If a unitary with this alphabet exists, then $\eta\le1$ and
$K\ge1$. The constants below may depend on $d,\Xi$.
The alphabet condition concerns the entries of the full physical
matrix in the computational basis. It does not concern the entries
of its local MPO tensors. A fixed local tensor alphabet alone
does not give the required positive lower bound for nonzero
physical entries; tensor contractions can produce arbitrarily
small nonzero entries.

\begin{theorem}
Let $U$ be a unitary on $N\ge1$ $d$-dimensional sites whose
consecutive operator Schmidt ranks are at most $D\ge1$.
Suppose $U_{x,y}\in\Xi\cup\{0\}$ for all computational basis
words $x,y$.
Then $U$ has an exact circuit with gate count and clean auxiliary
space polynomial jointly in $N,D$, including on a line with
only adjacent two-qudit gates. All auxiliaries are initialized
and returned to $|0\rangle$, and $CJ=JU$ retains the exact
scalar phase.

This is a nonuniform existence statement in the arbitrary-gate
model. It does not assert efficient computation of the required
canonical tensors from an arbitrary supplied MPO. It does not
assert the smaller $O_d(N\log(D+1))$ auxiliary bound.
The proof below is mathematical and is not Lean formalized.
\end{theorem}

\begin{proof}
Every row and column has at most $K$ nonzero entries, since
its squared norm is one and each nonzero entry has squared
absolute value at least $\eta^2$.

\subsection*{Exact row and column oracles}

Use right-canonical tensors for the normalized Choi coefficient
vector $\psi(x,y)=d^{-N/2}U_{x,y}$, with bond dimensions at most
$D$:
\[
 \psi(x,y)=A_1(x_1,y_1)\cdots A_N(x_N,y_N),\qquad
 \sum_{a,b}A_j(a,b)A_j(a,b)^\dagger=I.
\]
For a fixed column $y$ and an output prefix $z$, its prefix
mass is
\[
 p(z\mid y)=d^{-N}\sum_{x:\,x\text{ extends }z}|U_{x,y}|^2.
\]
It is either zero or at least $\eta^2d^{-N}$.
The conditioned transfer contraction and Hermitian rounding
estimate in
\nolinkurl{docs/audits/2026-10-03_mpu_joint_polynomial_permutations.tex}
give, with $L=c_d(D+1)^3$ and mesh $u=2^{-b}$,
\[
 e_{j-1}\le(1+Lu)e_j+Lu,\qquad
 e_0\le2NLu.
\]
Choosing $u\le\eta^2d^{-N}/(16NL)$ makes the numerical
prefix-mass error at most $\eta^2d^{-N}/8$.
A fixed rational threshold in
$[\eta^2d^{-N}/3,2\eta^2d^{-N}/3]$ therefore detects exactly
whether a prefix contains a nonzero entry. Such a threshold
can be represented with $O_{d,\Xi}(N)$ bits.
The comparison is an exact integer operation.

Explore the output-prefix tree in lexicographic order, retaining
only the positive prefixes. At any level there are at most $K$
such prefixes, since each contains a distinct nonzero leaf.
There are at most $dKN$ candidate prefix tests. A fixed
Boolean circuit realizes this search by using $K$ occupied
prefix slots at each level, testing their $d$ extensions and
compacting the positive results in order. Unoccupied slots
are explicitly marked and discarded. This yields the sorted
column support list, padded to $K$ slots with occupancy bits.

The exact alphabet labels are also recoverable. Scalar
contraction of the canonical matrices has a uniform error
$c_dN(D+1)^2u$. Approximate the fixed rescaling
$d^{N/2}$ with $O_{d,\Xi}(N)$ bits and rescale that scalar.
Increasing $b$ by $O_{d,\Xi}(\log(ND+1)+1)$ gives a value
within $\delta/8$ of the physical entry.
Approximate each member of $\Xi\cup\{0\}$ to the same
accuracy and compute the nearest label. The separation
$\delta$ makes that label unique and correct.
All comparisons, including squared complex distances, may
be made on finite integers with
\[
 B=O_{d,\Xi}(N+\log(N+1)+\log(D+1)+1)
\]
bits; a sufficiently accurate rational approximation to the
fixed alphabet is used in the circuit description.
Thus both the support and the alphabet label are exact
outputs of a Boolean calculation, although numerical
intermediates are rounded.

The same construction on $U^\dagger$ recovers sorted row
support and labels. Its canonical tensors are
$\overline{A_j(b,a)}$ and satisfy the same norm identities.
The recovered labels are conjugated back to obtain the
entries of the row of $U$; its support is unchanged.
No reversal of the virtual tensor order is needed.
One oracle evaluation has Boolean size at most
\[
 \operatorname{poly}_{d,\Xi}(N,D).
\]
For example, the direct double-sum transfer evaluation gives
$O_{d,\Xi}(KN^2(D+1)^4B^2)$ operations, with an increased
constant covering the list compaction and scalar label tests.
Retaining Boolean intermediate values, adding the result
to an output register, and reversing the calculation gives
exact reversible oracles with polynomial gate count and
clean polynomial workspace.

\subsection*{Two explicitly prepared isometries}

Let $X,Y$ be two $N$-qudit word registers and let a constant-size
tag register contain the three orthogonal states $0,1,2$.
Use an additional constant-size index register, initialized
at zero, whose dimension is at least $K+1$.
Its dimension and the tag dimension are fixed by $d,\Xi$,
independent of $N,D$.

Define two maps into the common output space by
\[
 S|y\rangle=
 \frac1{\sqrt K}\sum_x U_{x,y}|x,y,0,0\rangle
       +\sqrt{1-\frac1K}\,|0,y,1,0\rangle,
\]
\[
 T|x\rangle=
 \frac1{\sqrt K}\sum_{y:\,U_{x,y}\ne0}|x,y,0,0\rangle
       +\sqrt{1-\frac{k_x}{K}}\,|x,0,2,0\rangle,
\]
where $k_x$ is the row support size.
The fourth register is the index, reset to zero in these
specified output columns. Distinct domain words are
distinguished by the preserved word register. Column
normalization of $U$ and the definition of $k_x$ give
$S^\dagger S=T^\dagger T=I$. The different failure tags
give no cross terms. Hence
\[
 T^\dagger S=\frac1K U
\]
with the original complex phases.

These maps have polynomial clean unitary preparations from
the same initialized input inclusion
$J|y\rangle=|y,0,0,0\rangle$.
For $S$, first swap the two word registers, leaving the
column word $y$ in the second register. Reversibly compute
its sorted support and labels. If the nonzero labels are
$\alpha_0,\ldots,\alpha_{k_y-1}$, prepare the index state
\[
 \frac1{\sqrt K}\sum_{i<k_y}\alpha_i|i\rangle
       +\sqrt{1-\frac1K}|K\rangle.
\]
The squared norm is one because the column of $U$ has
norm one. There are only finitely many possible label
tuples, since $K,\Xi$ are fixed. Choose a full unitary
preparation for each valid tuple, and identity for invalid
tuples. The corresponding controlled operation acts on
only constantly many qudits and has constant gate cost.
It includes the exact phases of the $\alpha_i$.

On index $i<k_y$, write the $i$th support word into $X$
and retain tag zero. On index $K$, leave $X$ at zero
and write tag one. These are reversible controlled
additions. Erase the index by subtracting its address,
which is determined by the preserved column word, the
output word, and the tag: a successful output has a
unique sorted support address, and the failure tag has
address $K$. Address lookup is an exact Boolean
calculation with polynomial cost. Define it arbitrarily,
for example as zero, on unsupported word--tag pairs.
Reverse the support and label oracle using the preserved
column word. All its arithmetic work returns to zero.
On the initialized columns this gives exactly $S$.

For $T$, preserve the row word $x$ in $X$. Compute its
row support and prepare
\[
 \frac1{\sqrt K}\sum_{i<k_x}|i\rangle
       +\sqrt{1-\frac{k_x}{K}}|K\rangle.
\]
Write its support words into $Y$, use tag two for failure,
erase the index by the corresponding unique address,
and reverse the row oracle. The preparations depend on
only finitely many support counts and therefore have
constant cost apart from the polynomial oracles.

There is an important distinction between logical indices
and arithmetic work. Index erasure has been proved on
the specified initialized columns. It is not asserted
for arbitrary nonzero logical word, tag, or index inputs.
The index is therefore part of the logical register of
each full unitary preparation. Every core operation just
described is nevertheless a unitary on that whole register.
The separate arithmetic work returns to zero for every
logical input: the row or column word controlling its
oracle is preserved, and the support and label lists are
used only as controls. This remains true for arbitrary
coherent inputs and reference entanglement.

Thus there exist full logical unitaries $Z_S,Z_T$,
implemented with polynomial clean arithmetic workspace,
such that
\[
 Z_SJ=S,\qquad Z_TJ=T.
\]
Their inverses have the same clean workspace property.

\subsection*{Exact amplification}

Set $W=Z_T^\dagger Z_S$, $V=WJ$, and $P=JJ^\dagger$.
Then $V$ is an isometry, $P$ is an orthogonal projection,
and the preceding overlap gives
\[
 J^\dagger WJ=U/K,\qquad
 V^\dagger PV=K^{-2}I,\qquad
 KPV=JU.
\]
Apply the exact uniform-amplitude construction of Section~5 of
\nolinkurl{docs/audits/2026-10-02_mpu_rank_two_circuits.tex}, with
\[
 \ell=\lceil\pi K/4\rceil\le K,\qquad
 \theta=\frac{\pi}{4\ell+2},\qquad t=K\sin\theta\le1.
\]
One additional initialized attenuation flag gives success
amplitude $\sin\theta$. After $\ell$ reflection rounds,
the output is exactly the first-flag branch $JU$.
The exact two-plane identity retains the scalar phase;
it gives coefficient $1$, rather than an unspecified
unit complex scalar.

The successful projection $P$ tests all logical auxiliary
registers for zero. The prepared-image reflection is
implemented by $W$, the same initialized-register
reflection, and $W^\dagger$, together with the constant-size
attenuation operation. Such zero-register reflections
have polynomial gate count and polynomial clean
workspace by reversible conjunction. The arithmetic
workspace may be shared, since it erases for every
logical input after each invocation.
Since $K$ depends only on the fixed alphabet, the total
number of preparations and inverses is constant.
Their gate and workspace costs are polynomial in $N,D$.
At the end every logical auxiliary, index, tag, attenuation
flag, and arithmetic work register is zero.

Routing all two-qudit gates through the polynomial-size
register multiplies the gate count by at most another
polynomial factor. This proves the line version as well.
\end{proof}

The construction applies to a proper subclass of general
MPUs: a fixed physical-entry alphabet implies uniformly
bounded row and column support. It includes permutation
unitaries and monomial unitaries with a fixed phase
alphabet, and it can include coherent mixing among a
fixed number of basis outputs. It does not settle joint
polynomial synthesis for arbitrary dense MPUs, nor for
arbitrary sparse MPUs whose nonzero entries have no
fixed positive lower bound and no discrete label gap.

\begin{thebibliography}{2}
\bibitem{permutation-circuits}
TNLean contributors,
\emph{Joint polynomial circuits for permutation and monomial unitaries},
3 October 2026.
Local source:
\nolinkurl{docs/audits/2026-10-03_mpu_joint_polynomial_permutations.tex}.
\bibitem{exact-mpu-circuits}
TNLean contributors,
\emph{Exact polynomial circuits for matrix product unitaries},
2 October 2026, Section~5.
Local source:
\nolinkurl{docs/audits/2026-10-02_mpu_rank_two_circuits.tex}.
\end{thebibliography}
\end{document}
